\section{Khai báo sử dụng Generative AI}

Phần này phải phản ánh việc sử dụng AI thực tế. Không tạo prompt giả hoặc mô tả AI như tác giả thay cho thành viên. Mỗi nội dung cần cho thấy vấn đề ban đầu, công cụ được dùng, phạm vi hỗ trợ, phần nhóm tự kiểm chứng/chỉnh sửa và mức đóng góp cuối cùng.

\subsection{Công cụ, phạm vi và mức đóng góp}

\begin{longtable}{|L{25mm}|L{33mm}|L{40mm}|L{38mm}|}
  \caption{Tổng hợp việc sử dụng Generative AI}\\
  \hline
  \rowcolor{gray!15}
  \textbf{Công cụ} & \textbf{Phạm vi hỗ trợ} & \textbf{Đầu ra được sử dụng và mức đóng góp} & \textbf{Cách nhóm kiểm chứng/chỉnh sửa}\\\hline
  \endfirsthead
  \hline
  \rowcolor{gray!15}
  \textbf{Công cụ} & \textbf{Phạm vi hỗ trợ} & \textbf{Đầu ra được sử dụng và mức đóng góp} & \textbf{Cách nhóm kiểm chứng/chỉnh sửa}\\\hline
  \endhead
  Tên công cụ và phiên bản nếu biết & Phân tích, phản biện, code, diagram hoặc soát lỗi & Không dùng / tham khảo / chỉnh sửa nhiều / sử dụng có kiểm soát & Nguồn đối chiếu, test, review và quyết định cuối\\\hline
\end{longtable}

\subsection{Các phiên làm việc tiêu biểu}

Chọn một số phiên thực tế có giá trị: yêu cầu đầu vào, phần trả lời hữu ích, điểm sai hoặc chưa phù hợp, và phần nhóm phát triển thêm. Nếu đính kèm prompt, giữ đúng nội dung đã dùng và loại bỏ dữ liệu nhạy cảm.

\subsection{Trách nhiệm của nhóm}

Khẳng định thành viên chịu trách nhiệm về requirement, mô hình, thiết kế, code, test và khả năng giải thích; AI không thay thế việc kiểm chứng với đề bài, giảng viên, tài liệu kỹ thuật và kết quả chạy thực tế.
